\chapter{Chapter09 — Matching Asymptotics}
%%% generated by the blueprint pipeline from FABL.Chapter09.MatchingAsymptotics
\section{Overview}

Declarations of the Lean module \texttt{FABL.Chapter09.MatchingAsymptotics}.

\section{Declarations}

\begin{definition}[Fixed-shift factorial ratio]
\lean{FABL.fixedShiftFactorialRatio}
\label{FABL.fixedShiftFactorialRatio}
\leanok
For natural numbers $m$ and $d$, define the fixed-shift factorial ratio to be
$\displaystyle\prod_{i=0}^{3d-1}\frac{m+i+1}{m+3d+i+1}$, where an empty product is $1$.
\end{definition}

\begin{theorem}[Limit of a fixed-shift factorial ratio]
\lean{FABL.tendsto_fixedShiftFactorialRatio}
\label{FABL.tendsto_fixedShiftFactorialRatio}
\leanok
\uses{FABL.fixedShiftFactorialRatio}
\difficulty{5}
For every natural number $d$, the finite product
$\displaystyle\prod_{i=0}^{3d-1}\frac{m+i+1}{m+3d+i+1}$ converges to $1$ as the natural
number $m$ tends to infinity.
\end{theorem}

\begin{theorem}[Factorial expression for a finite product]
\lean{FABL.fixedShiftFactorialRatio_eq_factorial_ratio}
\label{FABL.fixedShiftFactorialRatio_eq_factorial_ratio}
\leanok
\uses{FABL.fixedShiftFactorialRatio}
\difficulty{5}
For natural numbers $m$ and $d$, the product of $(m+i+1)/(m+3d+i+1)$ over
$i=0,\ldots,3d-1$ equals $((m+3d)!)^2/((m+6d)!\,m!)$.
\end{theorem}

\begin{theorem}[Factorial identity for the six-block coefficient]
\lean{FABL.sixBlockMultinomial_shift_mul_factorials}
\label{FABL.sixBlockMultinomial_shift_mul_factorials}
\leanok
\uses{FABL.sixBlockMultinomial}
\difficulty{4}
For all nonnegative integers $m$ and $d$, the multinomial coefficient with six parts of
size $d$ and a seventh part of size $m$, multiplied by $(d!)^6m!$, equals $(m+6d)!$.
\end{theorem}

\begin{theorem}[Factorial identity for a binomial coefficient]
\lean{FABL.choose_three_mul_shift_mul_factorials}
\label{FABL.choose_three_mul_shift_mul_factorials}
\leanok
\difficulty{2}
For all natural numbers $m$ and $d$, $\binom{m+6d}{3d}(3d)!(m+3d)!=(m+6d)!$.
\end{theorem}

\begin{theorem}[Shifted factorial ratio identity]
\lean{FABL.exercise9_3_ratio_shift_eq}
\label{FABL.exercise9_3_ratio_shift_eq}
\leanok
\uses{FABL.choose_three_mul_shift_mul_factorials, FABL.fixedShiftFactorialRatio, FABL.fixedShiftFactorialRatio_eq_factorial_ratio, FABL.sixBlockMultinomial, FABL.sixBlockMultinomial_shift_mul_factorials}
\difficulty{5}
For nonnegative integers $m,d$, let
\[
R(m,d)=\prod_{i=0}^{3d-1}\frac{m+i+1}{m+3d+i+1},
\]
and let $M(m+6d,d)$ be the multinomial coefficient with six entries equal to $d$ and a
seventh equal to $m$. Then
\[
\frac{M(m+6d,d)}{\binom{m+6d}{3d}^{2}}=\frac{((3d)!)^2}{(d!)^6}R(m,d).
\]
\end{theorem}

\begin{theorem}[Six-block multinomial ratio limit]
\lean{FABL.tendsto_exercise9_3_multinomial_ratio}
\label{FABL.tendsto_exercise9_3_multinomial_ratio}
\leanok
\uses{FABL.exercise9_3_ratio_shift_eq, FABL.sixBlockMultinomial, FABL.tendsto_fixedShiftFactorialRatio}
\difficulty{3}
For every natural number $d$, as $m\to\infty$ the ratio of the multinomial coefficient
with six parts of size $d$ and a seventh part of size $m$ to $\binom{m+6d}{3d}^2$
converges to $\frac{((3d)!)^2}{(d!)^6}$.
\end{theorem}

\begin{definition}[Stirling approximation to the factorial]
\lean{FABL.factorialStirlingApprox}
\label{FABL.factorialStirlingApprox}
\leanok
For a natural number $n$, define its Stirling approximation to the factorial to be
$\sqrt{2n\pi}\,(n/e)^n$, where $e=\exp(1)$.
\end{definition}

\begin{theorem}[Asymptotic factorial formula]
\lean{FABL.factorial_isEquivalent_factorialStirlingApprox}
\label{FABL.factorial_isEquivalent_factorialStirlingApprox}
\leanok
\uses{FABL.factorialStirlingApprox}
\difficulty{1}
As $n$ tends to infinity through the natural numbers, $n!$ is asymptotic to
$\sqrt{2n\pi}\,(n/e)^n$; equivalently, the ratio of these quantities tends to $1$.
\end{theorem}

\begin{theorem}[Stirling approximation ratio]
\lean{FABL.factorialStirlingApprox_three_ratio_eq}
\label{FABL.factorialStirlingApprox_three_ratio_eq}
\leanok
\uses{FABL.factorialStirlingApprox}
\difficulty{5}
For each positive integer $d$, define $S(n)=\sqrt{2n\pi}\,(n/e)^n$. Then
$S(3d)^2/S(d)^6=\frac{3}{4\pi^2}\,\frac{3^{6d}}{d^2}$.
\end{theorem}

\begin{definition}[Level-three factorial constant]
\lean{FABL.levelThreeFactorialConstant}
\label{FABL.levelThreeFactorialConstant}
\leanok
For a natural number $d$, define the level-three factorial constant to be
$\displaystyle\frac{((3d)!)^2}{(d!)^6}$.
\end{definition}

\begin{definition}[Level-three growth quantity]
\lean{FABL.levelThreeGrowth}
\label{FABL.levelThreeGrowth}
\leanok
For a natural number $d$, define the level-three growth quantity to be
$\displaystyle\frac{3^{6d}}{d^2}$, with the quotient interpreted in the real numbers (so
its value is $0$ when $d=0$).
\end{definition}

\begin{theorem}[Factorial constant asymptotic equivalent]
\lean{FABL.levelThreeFactorialConstant_isEquivalent}
\label{FABL.levelThreeFactorialConstant_isEquivalent}
\leanok
\uses{FABL.factorialStirlingApprox, FABL.factorialStirlingApprox_three_ratio_eq, FABL.factorial_isEquivalent_factorialStirlingApprox, FABL.levelThreeFactorialConstant, FABL.levelThreeGrowth}
\difficulty{4}
As $d$ tends to infinity through the natural numbers, $\frac{(3d)!^2}{(d!)^6}$ is
asymptotic to $\frac{3}{4\pi^2}\,\frac{3^{6d}}{d^2}$, meaning that the ratio of the two
quantities tends to $1$.
\end{theorem}

\begin{theorem}[Factorial constant growth order]
\lean{FABL.levelThreeFactorialConstant_isTheta}
\label{FABL.levelThreeFactorialConstant_isTheta}
\leanok
\uses{FABL.levelThreeFactorialConstant, FABL.levelThreeFactorialConstant_isEquivalent, FABL.levelThreeGrowth}
\difficulty{3}
As $d$ tends to infinity through the natural numbers, $\frac{(3d)!^2}{(d!)^6}$ is of the
same asymptotic order as $\frac{3^{6d}}{d^2}$; their ratio is eventually bounded above
and below by positive constants.
\end{theorem}

\begin{definition}[Level-three book growth quantity]
\lean{FABL.levelThreeBookGrowth}
\label{FABL.levelThreeBookGrowth}
\leanok
For a natural number $d$, define the level-three book growth quantity to be
$\displaystyle\frac{9^{3d}}{(3d)^2}$, with the quotient interpreted in the real numbers
(so its value is $0$ when $d=0$).
\end{definition}

\begin{theorem}[Level three growth identity]
\lean{FABL.levelThreeGrowth_eq_nine_mul_bookGrowth}
\label{FABL.levelThreeGrowth_eq_nine_mul_bookGrowth}
\leanok
\uses{FABL.levelThreeBookGrowth, FABL.levelThreeGrowth}
\difficulty{4}
For every natural number $d$, define $G(d)=d^{-2}3^{6d}$ and $B(d)=(3d)^{-2}9^{3d}$,
with the quotients interpreted as real expressions. Then $G(d)=9B(d)$.
\end{theorem}

\begin{theorem}[Factorial constant and book growth]
\lean{FABL.levelThreeFactorialConstant_isTheta_bookGrowth}
\label{FABL.levelThreeFactorialConstant_isTheta_bookGrowth}
\leanok
\uses{FABL.levelThreeBookGrowth, FABL.levelThreeFactorialConstant, FABL.levelThreeFactorialConstant_isTheta, FABL.levelThreeGrowth, FABL.levelThreeGrowth_eq_nine_mul_bookGrowth}
\difficulty{4}
As $d$ tends to infinity through the natural numbers, $\frac{(3d)!^2}{(d!)^6}$ is of the
same asymptotic order as $\frac{9^{3d}}{(3d)^2}$; their ratio is eventually bounded
above and below by positive constants.
\end{theorem}

\begin{theorem}[Uniform norm lower bound]
\lean{FABL.exercise9_3_uniformLpNorm_lower_bound}
\label{FABL.exercise9_3_uniformLpNorm_lower_bound}
\leanok
\uses{FABL.exercise9_3_fourthMoment_lower_bound, FABL.levelKPolynomial, FABL.sixBlockMultinomial, FABL.uniformLpNorm, FABL.uniformLpNorm_even_pow_eq_expect_even, FABL.uniformLpNorm_nonneg, FABL.uniformLpNorm_two_sq_eq_expect_sq}
\difficulty{4}
Let $n,d$ be nonnegative integers with $6d\leq n$, and let $F$ be the level-$3d$
polynomial on the discrete cube of dimension $n$. If $M(n,d)$ denotes the multinomial
coefficient with six entries equal to $d$ and a seventh equal to $n-6d$, then
\[
\left(\frac{M(n,d)}{\binom{n}{3d}^{2}}\right)^{1/4}\|F\|_2\leq \|F\|_4,
\]
where the norms are uniform norms on the cube.
\end{theorem}

\begin{definition}[Level-three book norm growth quantity]
\lean{FABL.levelThreeBookNormGrowth}
\label{FABL.levelThreeBookNormGrowth}
\leanok
For a natural number $d$, define the level-three book norm growth quantity to be
$\displaystyle(3d)^{-1/2}(\sqrt{3})^{3d}$, where the first factor is the real power
$\operatorname{rpow}(3d,-1/2)$ (and equals $0$ when $d=0$).
\end{definition}

\begin{theorem}[Fourth root of the level-three growth]
\lean{FABL.levelThreeBookGrowth_fourthRoot_eq_normGrowth}
\label{FABL.levelThreeBookGrowth_fourthRoot_eq_normGrowth}
\leanok
\uses{FABL.levelThreeBookGrowth, FABL.levelThreeBookNormGrowth}
\difficulty{4}
For every positive integer $d$, the fourth root, interpreted as a real power, of
$\frac{1}{(3d)^2}9^{3d}$ equals $(3d)^{-1/2}(\sqrt{3})^{3d}$.
\end{theorem}

\begin{theorem}[Factorial growth and norm growth]
\lean{FABL.levelThreeFactorialConstant_fourthRoot_isTheta_normGrowth}
\label{FABL.levelThreeFactorialConstant_fourthRoot_isTheta_normGrowth}
\leanok
\uses{FABL.levelThreeBookGrowth, FABL.levelThreeBookGrowth_fourthRoot_eq_normGrowth, FABL.levelThreeBookNormGrowth, FABL.levelThreeFactorialConstant, FABL.levelThreeFactorialConstant_isTheta_bookGrowth}
\difficulty{4}
As $d$ tends to infinity through the natural numbers, the function
$\left(\frac{(3d)!^2}{(d!)^6}\right)^{1/4}$ is of the same asymptotic order as
$(3d)^{-1/2}(\sqrt{3})^{3d}$; that is, their ratio is eventually bounded above and below
by positive constants.
\end{theorem}

\begin{definition}[Matching factorial ratio]
\lean{FABL.matchingFactorialRatio}
\label{FABL.matchingFactorialRatio}
\leanok
For natural numbers $r$ and $k$, define the matching factorial ratio to be
$\displaystyle\frac{(rk)!}{(k!)^r r^{rk}}$, with the quotient taken in $\mathbb{R}$.
\end{definition}

\begin{theorem}[Improvement constant as a power]
\lean{FABL.matchingImprovementConstant_eq_rpow_factorialRatio}
\label{FABL.matchingImprovementConstant_eq_rpow_factorialRatio}
\leanok
\uses{FABL.matchingFactorialRatio, FABL.matchingImprovementConstant}
\difficulty{1}
For natural numbers $r$ and $k$, the improvement constant is equal to
$\left(\frac{(rk)!}{(k!)^r r^{rk}}\right)^{1/(2r)}$, with powers interpreted as real
powers.
\end{theorem}

\begin{definition}[Matching Stirling core]
\lean{FABL.matchingStirlingCore}
\label{FABL.matchingStirlingCore}
\leanok
For natural numbers $r$ and $k$, define the matching Stirling core by
$\sqrt{2rk\pi}/(\sqrt{2k\pi})^r$.
\end{definition}

\begin{theorem}[Factorial approximation ratio identity]
\lean{FABL.factorialStirlingApprox_matching_ratio_eq_core}
\label{FABL.factorialStirlingApprox_matching_ratio_eq_core}
\leanok
\uses{FABL.factorialStirlingApprox, FABL.matchingStirlingCore}
\difficulty{4}
For each nonnegative integer $m$, write $A(m)=\sqrt{2m\pi}\,(m/e)^m$, where $e=\exp(1)$.
For positive integers $r$ and $k$, define $C(r,k)=\sqrt{2rk\pi}/(\sqrt{2k\pi})^r$. Then
$A(rk)/(A(k)^r r^{rk})=C(r,k)$.
\end{theorem}

\begin{theorem}[Factorial ratio asymptotic equivalent]
\lean{FABL.matchingFactorialRatio_isEquivalent_core}
\label{FABL.matchingFactorialRatio_isEquivalent_core}
\leanok
\uses{FABL.factorialStirlingApprox, FABL.factorialStirlingApprox_matching_ratio_eq_core, FABL.factorial_isEquivalent_factorialStirlingApprox, FABL.matchingFactorialRatio, FABL.matchingStirlingCore}
\difficulty{4}
Fix a positive natural number $r$. As the natural number $k$ tends to infinity, the
ratio $\frac{(rk)!}{(k!)^r r^{rk}}$ is asymptotic to
$\frac{\sqrt{2rk\pi}}{(\sqrt{2k\pi})^r}$; that is, the quotient of these two
quantities tends to $1$.
\end{theorem}

\begin{definition}[Matching core growth]
\lean{FABL.matchingCoreGrowth}
\label{FABL.matchingCoreGrowth}
\leanok
For natural numbers $r$ and $k$, define the matching core growth $G(r,k)$ by
$G(r,k)=k^{(1-r)/2}$, where the power is taken in the real numbers.
\end{definition}

\begin{theorem}[Exact matching core factorization]
\lean{FABL.matchingStirlingCore_eq_constant_mul_growth}
\label{FABL.matchingStirlingCore_eq_constant_mul_growth}
\leanok
\uses{FABL.matchingCoreGrowth, FABL.matchingStirlingCore}
\difficulty{5}
For positive integers $r$ and $k$, the matching Stirling core equals $\sqrt
r\,(2\pi)^{(1-r)/2}$ times the matching core growth.
\end{theorem}

\begin{theorem}[Asymptotic matching core growth]
\lean{FABL.matchingStirlingCore_isTheta_growth}
\label{FABL.matchingStirlingCore_isTheta_growth}
\leanok
\uses{FABL.matchingCoreGrowth, FABL.matchingStirlingCore, FABL.matchingStirlingCore_eq_constant_mul_growth}
\difficulty{4}
For every positive integer $r$, as the integer $k$ tends to infinity, the matching
Stirling core is asymptotically comparable to the matching core growth; that is, their
ratio is bounded above and below by positive constants for all sufficiently large $k$.
\end{theorem}

\begin{theorem}[Factorial ratio growth order]
\lean{FABL.matchingFactorialRatio_isTheta_growth}
\label{FABL.matchingFactorialRatio_isTheta_growth}
\leanok
\uses{FABL.matchingCoreGrowth, FABL.matchingFactorialRatio, FABL.matchingFactorialRatio_isEquivalent_core, FABL.matchingStirlingCore_isTheta_growth}
\difficulty{1}
Fix a positive natural number $r$. As $k$ tends to infinity through the natural numbers,
the function $\frac{(rk)!}{(k!)^r r^{rk}}$ has the same order of growth, up to positive
constant factors, as $k^{(1-r)/2}$.
\end{theorem}

\begin{definition}[Matching improvement growth]
\lean{FABL.matchingImprovementGrowth}
\label{FABL.matchingImprovementGrowth}
\leanok
For natural numbers $r$ and $k$, define the matching improvement growth to be
$k^{-1/4+1/(4r)}$, with the exponent interpreted using real arithmetic and the power as
a real power.
\end{definition}

\begin{theorem}[Power identity for matching growth]
\lean{FABL.matchingCoreGrowth_rpow_eq_improvementGrowth}
\label{FABL.matchingCoreGrowth_rpow_eq_improvementGrowth}
\leanok
\uses{FABL.matchingCoreGrowth, FABL.matchingImprovementGrowth}
\difficulty{3}
For natural numbers $r$ and $k$ with $r>0$, one has
$\left(k^{(1-r)/2}\right)^{1/(2r)}=k^{-1/4+1/(4r)}$, where powers are real powers.
\end{theorem}

\begin{theorem}[Improvement constant growth order]
\lean{FABL.matchingImprovementConstant_isTheta}
\label{FABL.matchingImprovementConstant_isTheta}
\leanok
\uses{FABL.matchingCoreGrowth, FABL.matchingCoreGrowth_rpow_eq_improvementGrowth, FABL.matchingFactorialRatio, FABL.matchingFactorialRatio_isTheta_growth, FABL.matchingImprovementConstant, FABL.matchingImprovementConstant_eq_rpow_factorialRatio, FABL.matchingImprovementGrowth}
\difficulty{5}
Fix a positive natural number $r$. As $k$ tends to infinity through the natural numbers,
the improvement constant has the same order of growth, up to positive constant factors,
as $k^{-1/4+1/(4r)}$.
\end{theorem}
